\documentclass[12pt, a4paper, oneside]{article}
\usepackage{amsmath, amsthm, amssymb, bm, ctex}
\usepackage[indent=0pt]{parskip}

\newcounter{probid}
\newcounter{subprobid}[probid]

\author{杨濛 2021213247}
\title{数学作业\\第二周}
\date{\today}

\newenvironment{problem}{
    \stepcounter{probid}
    \noindent\textbf{\arabic{probid}. }
}{\par}
\newenvironment{subproblem}{
    \stepcounter{subprobid}
    (\arabic{subprobid})
}{\par}
\newenvironment{solution}{
    \noindent\textbf{解: }
}{\par}

\begin{document}
    \maketitle
    \begin{problem}
        用检根法确定下列级数的收敛性：
    \end{problem}
    \begin{subproblem}
        $\sum\limits_{n=1}^\infty \dfrac{n}{2^n}$
    \end{subproblem}
    \begin{solution}
        \[\lim\limits_{n\to\infty} \sqrt[n]{\dfrac{2}{2^n}}
        = \dfrac{1}{2}\lim\limits_{n\to\infty}\sqrt[n]{n} = \dfrac{1}{2}<1\]
        所以原级数收敛.
    \end{solution}
    \begin{subproblem}
        \(\sum\limits_{n=1}^\infty\left( \dfrac{n}{3n-1} \right)^{2n-1}\)
    \end{subproblem}
    \begin{solution}
        \[\lim\limits_{n\to\infty}
        \sqrt[n]{\left( \dfrac{n}{3n-1} \right)^{2n-1}}
        = \lim\limits_{n\to\infty}
        \left( \dfrac{n}{3n-1} \right)^{\frac{2n-1}{n}}
        = \dfrac{1}{9}<1\]
        所以原级数收敛.
    \end{solution}
    \begin{problem}
        利用级数收敛的必要条件,证明:\(\lim\limits_{n\to\infty}
        \dfrac{n^n}{(n!)^2}= 0\)
    \end{problem}
    \begin{solution}
        令\(a_n = \dfrac{n^n}{(n!)^2}\),则
        \begin{align*}
            \lim\limits_{n\to\infty}\dfrac{a_{n+1}}{a_n}
            & = \lim\limits_{n\to\infty}
            \dfrac{\left(n+1\right)^{n+1}}
            {\left[\left(n+1\right)!\right]^2}
            \cdot \dfrac{\left(n!\right)^2}{n^n}\\
            & = \lim\limits_{n\to\infty}
            \dfrac{\left(n+1\right)^{n-1}}{n^n}\\
            & = \lim\limits_{n\to\infty}
            \dfrac{\left(1+\frac{1}{n}\right)^{n-1}}{n}\\
            & = 0
        \end{align*}
        根据检比法,级数\(\sum\limits_{n=1}^\infty a_n\)收敛.

        根据级数收敛的必要条件, \(\lim\limits_{n\to\infty}
        \dfrac{n^n}{(n!)^2}= 0\).
    \end{solution}
    \begin{problem}
        设正项级数\(\sum\limits_{n=1}^\infty u_n\)收敛,证明:当$p>\dfrac{1}{2}$时,
        级数$\sum\limits_{n=1}^\infty \dfrac{\sqrt{u_n}}{n^p}$收敛.
    \end{problem}
    \begin{solution}
        由于$p>\dfrac{1}{2}$,故对于所有$u_n$,有
        \[\sum\limits_{n=1}^\infty\dfrac{u_n}{n^{2p}}
        < \sum\limits_{n=1}^\infty u_n\]
        由比较审敛法可知级数\(\sum\limits_{n=1}^\infty \dfrac{u_n}{n^{2p}}\)收敛.

        考虑级数$\sum\limits_{n=1}^\infty \dfrac{\sqrt{u_n}}{n^p}$的部分和$S(n)$,有
        \begin{align*}
            S(n) & = \sum\limits_{k=1}^n \dfrac{\sqrt{u_n}}{n^p} \\
            & = \sum\limits_{k=1}^n \sqrt{\dfrac{u_n}{n^{2p}}} \\
            & < \sqrt{\dfrac{\sum\limits_{k=1}^n \dfrac{u_n}{n^{2p}}}{n}}
        \end{align*}
    \end{solution}
    \begin{problem}
        判定以下级数是否收敛。如果收敛，请判断是绝对收敛还是条件收敛。
    \end{problem}
    \begin{subproblem}
        \(\sum\limits_{n=1}^\infty
        \left(-1\right)^{n-1}\dfrac{\left(2n\right)!!}{\left(2n-1\right)!!}\)
    \end{subproblem}
    \begin{solution}
        令\(u_n = \dfrac{\left(2n\right)!!}{\left(2n-1\right)!!}\),
        则\(\lim\limits_{n\to\infty} u_n > \dfrac{2}{1} = 2 \ne 0\),所以该级数发散.
    \end{solution}
    \begin{subproblem}
        \(\sum\limits_{n=1}^\infty \dfrac{\left(-1\right)^{n-1}}
        {\ln\left(n+1\right)}\)
    \end{subproblem}
    \begin{solution}
        由于\(\ln x\)在$x\in \mathbf{R}$上单调递增,所以\(\dfrac{1}{\ln x}\)
        在$x \in \mathbf{R}$上单调递减. 根据柯西判别法,原交错级数收敛.

        由\(\ln x < x\)得
        \[\sum\limits_{n=1}^\infty \dfrac{1}{\ln n}
        > \sum\limits_{n=1}^\infty \dfrac{1}{n}\]
        由调和级数发散以及比较准则,原级数绝对发散.
        
        综上,原级数条件收敛.
    \end{solution}
    \begin{subproblem}
        \(\sum\limits_{n=1}^\infty\dfrac{n!\times 2^n}{n^n}\sin\dfrac{n\pi}{3}\)
    \end{subproblem}
    \begin{solution}
        显然\(\dfrac{n!\times 2^n}{n^n}\sin\dfrac{n\pi}{3}
        < \dfrac{n!\times 2^n}{n^n}\).

        令\(u_n = \dfrac{n!\times 2^n}{n^n}\), 则
        \[\lim\limits_{n\to\infty} \dfrac{u_{n+1}}{u_n}
        = \lim\limits_{n\to\infty} \dfrac{2n^n}{(n+1)^n}
        = \dfrac{2}{\rm{e}} < 1\]
        故根据检比法,原级数绝对收敛.
    \end{solution}
    \begin{problem}
        设级数\(\sum\limits_{n=2}^\infty\left(u_n-u_{n-1}\right)\)收敛,
        正项级数\(\sum\limits_{n=1}^\infty v_n\)收敛,试证明
        级数\(\sum\limits_{n=1}^\infty u_n v_n\)绝对收敛.
    \end{problem}
    \begin{solution}
        由\(\sum\limits_{n=2}^\infty\left(u_n-u_{n-1}\right)\)收敛可知
        其部分和的极限
        \[\lim\limits_{n\to\infty}S(n)
        = \lim\limits_{n\to\infty}u_n-u_1= 0
        \Rightarrow \lim\limits_{n\to\infty} u_n = u_1\]
        此时根据数列极限的有界性:
        \[\forall n \in \mathbf{Z}_+, \exists M \in R+, 
        \left|u_n\right| < M\]
        因此
        \[\sum\limits_{n=1}^\infty \left|u_n v_n\right|
        < M\sum\limits_{n=1}^\infty \left|v_n\right|\]
        由比较准则可知原级数绝对收敛.
    \end{solution}
\end{document}